\documentclass[12pt,a4paper]{article}
\usepackage[utf8]{inputenc}
\usepackage[T5]{fontenc}
\usepackage[vietnamese]{babel}
\usepackage[margin=1.3cm,top=0.8cm,bottom=1.0cm,includehead,includefoot,headheight=18pt,headsep=6pt,footskip=16pt]{geometry}
\usepackage{amsmath,amssymb,array,tabularx,enumitem,fancyhdr,tikz,xcolor,needspace,multicol}
\usetikzlibrary{calc,angles,quotes,arrows.meta,babel,patterns}

\definecolor{ink}{RGB}{30,55,70}
\definecolor{accent}{RGB}{18,82,112}

\setlength{\parindent}{0pt}
\setlength{\parskip}{2pt}

\newcommand{\headerQuote}{%
  \ifcase\numexpr\value{page}-1\relax
    Math gives us hope that every problem has a solution%
  \or Mathematics is about understanding, not just numbers%
  \or The essence of mathematics lies in its freedom%
  \or Pure mathematics is the poetry of logical ideas%
  \or In mathematics, asking good questions matters most%
  \or Everything around you is built on mathematics%
  \or Mathematics is the queen of the sciences%
  \or Stay with tough problems a little bit longer%
  \or Success is the sum of small daily efforts%
  \or Think deeply, check carefully, solve cleanly%
  \or Genius is 1\% inspiration and 99\% perspiration%
  \or Practice makes perfect — keep solving!%
  \or Math + You = Champion%
  \or Where Frustration Becomes Inspiration%
  \or Turning Confusion into Clarity%
  \or Don't get mad. Learn to add%
  \or The only way to learn mathematics is to do mathematics%
  \or Small steps every day lead to big results%
  \or Every expert was once a beginner%
  \or Mathematics is the music of reason%
  \or Math gives us hope that every problem has a solution%
  \or Mathematics is about understanding, not just numbers%
  \or The essence of mathematics lies in its freedom%
  \or Pure mathematics is the poetry of logical ideas%
  \or In mathematics, asking good questions matters most%
  \or Everything around you is built on mathematics%
  \or Mathematics is the queen of the sciences%
  \or Stay with tough problems a little bit longer%
  \or Success is the sum of small daily efforts%
  \or Think deeply, check carefully, solve cleanly%
  \or Genius is 1\% inspiration and 99\% perspiration%
  \or Practice makes perfect — keep solving!%
  \or Math + You = Champion%
  \or Where Frustration Becomes Inspiration%
  \or Turning Confusion into Clarity%
  \or Don't get mad. Learn to add%
  \or The only way to learn mathematics is to do mathematics%
  \or Small steps every day lead to big results%
  \or Every expert was once a beginner%
  \or Mathematics is the music of reason%
  \else Practice makes perfect — keep solving!%
  \fi
}

\pagestyle{fancy}
\fancyhf{}
\lhead{\footnotesize\textit{\color{gray!70!black}\headerQuote}}
\rhead{\small\textbf{\color{ink}Người biên soạn: Dương Gia Bảo}}
\lfoot{}
\rfoot{\small Trang \thepage/20}
\renewcommand{\headrulewidth}{0.4pt}
\renewcommand{\footrulewidth}{0pt}

\fancypagestyle{dapan}{
  \fancyhf{}
  \lhead{\footnotesize\textit{\color{gray!70!black}\headerQuote}}
  \rhead{\small\textbf{\color{ink}Người biên soạn: Dương Gia Bảo}}
  \lfoot{\small\textit{\color{gray!60!black}Đáp án Đề cương Ôn tập Chương IV (Bài 7 – 10) — Toán 10 KNTT}}
  \rfoot{\small Trang \thepage/2}
  \renewcommand{\headrulewidth}{0.4pt}
  \renewcommand{\footrulewidth}{0pt}
}

\newcommand{\chapterbanner}[2]{%
  \par\vspace{3pt}\noindent
  \begin{tikzpicture}
    \node[draw=accent, fill=accent!6, rounded corners=4pt, inner xsep=10pt, inner ysep=5pt, text width=\dimexpr\linewidth-20.8pt\relax, line width=0.8pt] {%
      {\large\bfseries\color{ink}#1}\\[2pt]
      {\small\color{gray!75!black}#2}%
    };
  \end{tikzpicture}\par\nopagebreak\vspace{4pt}
}

\newcommand{\parttitle}[1]{\par\vspace{3pt}\needspace{3\baselineskip}{\normalsize\bfseries\color{accent}#1}\par\nopagebreak\vspace{2pt}}

\newcommand{\answerbox}{%
  {\unskip\nobreak\hfil\penalty50\hskip1.2em\hbox{}\nobreak\hfil
  \mbox{\textit{Đáp số:}\ \fbox{\rule{0pt}{8.5pt}\hspace{2.2cm}}}%
  \parfillskip=0pt\finalhyphendemerits=0\par}%
}

\setlist[enumerate,1]{label=\textbf{Câu \arabic*.},leftmargin=1.15cm,labelsep=.15cm,itemsep=2pt,topsep=2pt,parsep=0pt}

\newcommand{\oli}[1]{%
  \par\nopagebreak\vspace{2.5pt}\noindent
  \begin{tikzpicture}
    \draw[step=0.5cm, color=cyan!25!gray!25, line width=0.25pt] (0,0) grid (\linewidth, #1);
    \draw[color=cyan!50!gray!45, line width=0.45pt, rounded corners=2pt] (0,0) rectangle (\linewidth, #1);
    \node[anchor=north west, font=\scriptsize\itshape\color{gray!75}, inner sep=3pt] at (0, #1) {Bài làm:};
  \end{tikzpicture}\par\vspace{3.5pt}
}

\newcolumntype{C}[1]{>{\centering\arraybackslash}m{#1}}


\begin{document}

\begin{center}
  {\Large\bfseries\color{ink} ĐỀ CƯƠNG ÔN TẬP CHƯƠNG IV — VECTƠ (40 CÂU)}\\[3pt]
  {\small\color{gray!75!black} \textit{Hình thức: 40 câu trả lời ngắn (Bài 7, 8, 9, 10 — Sách Kết nối tri thức với cuộc sống)}}
\end{center}
\vspace{2pt}

\chapterbanner{ÔN TẬP KIẾN THỨC TRỌNG TÂM: BÀI 7 ĐẾN BÀI 10}{(Học sinh trình bày tóm tắt các bước giải vào khung ô li và điền đáp số cuối cùng vào ô tương ứng)}

\parttitle{CHỦ ĐỀ 1. CÁC KHÁI NIỆM MỞ ĐẦU (Bài 7 — Câu 1 đến Câu 10)}

\begin{enumerate}

\item Cho thập giác đều $A_1A_2\ldots A_{10}$ có $10$ đỉnh. Có bao nhiêu vectơ khác vectơ $\vec{0}$ có điểm đầu và điểm cuối là các đỉnh của thập giác đều đã cho?\answerbox
\oli{8.0cm}

\item Cho hình chữ nhật $ABCD$ có tâm $O$. Có bao nhiêu vectơ khác vectơ $\vec{0}$ cùng phương với vectơ $\vec{AB}$ mà có điểm đầu và điểm cuối là các điểm thuộc tập hợp $\{A, B, C, D, O\}$?\answerbox
\oli{8.0cm}

\newpage

\item Cho tam giác $ABC$. Gọi $M, N, P$ lần lượt là trung điểm của các cạnh $BC, CA, AB$. Có bao nhiêu vectơ khác vectơ $\vec{MN}$ bằng vectơ $\vec{MN}$ mà có điểm đầu và điểm cuối là các điểm thuộc tập hợp $\{A, B, C, M, N, P\}$?\answerbox
\oli{9.0cm}

\item \noindent\begin{minipage}[t]{0.62\linewidth}
Cho lưới ô vuông gồm các hình vuông đơn vị có cạnh bằng $1$ (như hình vẽ bên). Hai điểm $A, B$ là hai mắt lưới. Tính độ dài của vectơ $\vec{u} = \vec{AB}$.\answerbox
\end{minipage}\hfill
\begin{minipage}[t]{0.34\linewidth}\raggedleft
\begin{tikzpicture}[scale=0.55, >=Stealth]
  \draw[step=1, color=gray!35, line width=0.35pt] (0,0) grid (6,6);
  \draw[->, accent, line width=1.3pt] (1,1) -- (4,5);
  \fill[accent] (1,1) circle (2.5pt) node[below left, font=\scriptsize] {$A$};
  \fill[accent] (4,5) circle (2.5pt) node[above right, font=\scriptsize] {$B$};
  \node[above left, font=\scriptsize, accent] at (2.5,3.2) {$\vec{u}$};
  \draw[<->, gray!80!black, line width=0.5pt] (0,-0.4) -- (1,-0.4) node[midway, below, font=\scriptsize] {$1$};
\end{tikzpicture}
\end{minipage}
\oli{8.0cm}

\newpage

\item \noindent\begin{minipage}[t]{0.62\linewidth}
Cho hình thang vuông $ABCD$ vuông tại $A$ và $D$ có đáy nhỏ $AB = 3\text{ cm}$, đáy lớn $CD = 7\text{ cm}$ và cạnh bên $AD = 3\text{ cm}$ (như hình vẽ bên). Tính độ dài của vectơ $\vec{BC}$ (theo đơn vị xentimét).\answerbox
\end{minipage}\hfill
\begin{minipage}[t]{0.34\linewidth}\raggedleft
\begin{tikzpicture}[scale=0.55, >=Stealth]
  \coordinate (A) at (0,3);
  \coordinate (D) at (0,0);
  \coordinate (C) at (7,0);
  \coordinate (B) at (3,3);
  \draw[thick] (A) -- (B) -- (C) -- (D) -- cycle;
  \draw[dashed, gray] (B) -- (3,0) node[below, font=\scriptsize] {$H$};
  \draw[->, red!75!black, line width=1.2pt] (B) -- (C);
  \fill[ink] (A) circle (2pt) node[above left, font=\scriptsize] {$A$};
  \fill[ink] (B) circle (2pt) node[above right, font=\scriptsize] {$B$};
  \fill[ink] (C) circle (2pt) node[below right, font=\scriptsize] {$C$};
  \fill[ink] (D) circle (2pt) node[below left, font=\scriptsize] {$D$};
  \node[above, font=\scriptsize] at (1.5,3) {$3$};
  \node[left, font=\scriptsize] at (0,1.5) {$3$};
  \node[below, font=\scriptsize] at (3.5,0) {$7$};
\end{tikzpicture}
\end{minipage}
\oli{8.0cm}

\item Cho tam giác $ABC$ vuông tại $A$ có $AB = 9\text{ cm}$ và $AC = 12\text{ cm}$. Gọi $M$ là trung điểm của cạnh huyền $BC$. Tính độ dài của vectơ $\vec{AM}$ (theo đơn vị xentimét).\answerbox
\oli{9.0cm}

\newpage

\item Cho hình vuông $ABCD$ có cạnh bằng $8\text{ cm}$. Gọi $M$ là trung điểm của cạnh $AB$ và $N$ là trung điểm của cạnh $BC$. Tính bình phương độ dài của vectơ $\vec{MN}$ (tức tính $|\vec{MN}|^2$).\answerbox
\oli{9.0cm}

\item Cho tam giác $ABC$ có ba cạnh $AB = 5$, $AC = 7$ và $BC = 8$. Gọi $D, E$ lần lượt là trung điểm của các cạnh $AB$ và $AC$. Tính độ dài của vectơ $\vec{DE}$.\answerbox
\oli{9.0cm}

\newpage

\item Cho hình chữ nhật $ABCD$ có diện tích bằng $48\text{ cm}^2$ và cạnh $AB = 6\text{ cm}$. Tính độ dài của vectơ đường chéo $\vec{BD}$ (theo đơn vị xentimét).\answerbox
\oli{9.0cm}

\item Cho $6$ điểm phân biệt $A, B, C, D, E, F$ trong đó không có ba điểm nào thẳng hàng. Có bao nhiêu vectơ khác vectơ $\vec{0}$ có điểm đầu là điểm $A$ và điểm cuối là một trong các điểm còn lại?\answerbox
\oli{9.0cm}

\end{enumerate}

\newpage

\parttitle{CHỦ ĐỀ 2. TỔNG VÀ HIỆU CỦA HAI VECTƠ (Bài 8 — Câu 11 đến Câu 20)}

\begin{enumerate}[resume*]

\item Cho hình chữ nhật $ABCD$ có $AB = 5$ và $AD = 12$. Tính độ dài của vectơ hiệu $\vec{u} = \vec{AB} - \vec{CB}$.\answerbox
\oli{8.0cm}

\item \noindent\begin{minipage}[t]{0.68\linewidth}
Cho hình thoi $ABCD$ có cạnh bằng $10\text{ cm}$ và đường chéo $BD = 12\text{ cm}$ (như hình vẽ bên). Tính độ dài của vectơ tổng $\vec{w} = \vec{AB} + \vec{AD}$ (theo đơn vị xentimét).\answerbox
\end{minipage}\hfill
\begin{minipage}[t]{0.28\linewidth}\raggedleft
\begin{tikzpicture}[scale=0.32, >=Stealth]
  \coordinate (A) at (0,8);
  \coordinate (B) at (6,0);
  \coordinate (C) at (0,-8);
  \coordinate (D) at (-6,0);
  \coordinate (O) at (0,0);
  \draw[thick] (A) -- (B) -- (C) -- (D) -- cycle;
  \draw[dashed, gray] (A) -- (C);
  \draw[dashed, gray] (D) -- (B);
  \fill[ink] (A) circle (2.5pt) node[above, font=\scriptsize] {$A$};
  \fill[ink] (B) circle (2.5pt) node[right, font=\scriptsize] {$B$};
  \fill[ink] (C) circle (2.5pt) node[below, font=\scriptsize] {$C$};
  \fill[ink] (D) circle (2.5pt) node[left, font=\scriptsize] {$D$};
  \fill[ink] (O) circle (2pt) node[above right, font=\scriptsize] {$O$};
  \node[above right, font=\scriptsize] at (3,4) {$10$};
  \node[below, font=\scriptsize] at (3,0) {$6$};
\end{tikzpicture}
\end{minipage}
\oli{7.5cm}

\newpage

\item Cho bốn điểm phân biệt $A, B, C, D$ thỏa mãn đoạn thẳng $BC$ có độ dài bằng $9\text{ cm}$. Tính độ dài của vectơ tổng $\vec{x} = \vec{AB} + \vec{CD} - \vec{AD}$ (theo đơn vị xentimét).\answerbox
\oli{9.0cm}

\item Cho hình chữ nhật $ABCD$ có $AB = 8$ và $AD = 6$. Tính độ dài của vectơ hiệu $\vec{y} = \vec{AB} - \vec{AD}$.\answerbox
\oli{9.0cm}

\newpage

\item Hai lực đồng quy $\vec{F}_1$ và $\vec{F}_2$ cùng tác dụng vào một chất điểm tại điểm $O$, có phương vuông góc với nhau và có độ lớn lần lượt là $F_1 = 30\text{ N}$ và $F_2 = 40\text{ N}$. Tính độ lớn của hợp lực $\vec{F} = \vec{F}_1 + \vec{F}_2$ (theo đơn vị Niutơn).\answerbox
\oli{9.0cm}

\item Hai lực $\vec{F}_1$ và $\vec{F}_2$ có cùng độ lớn bằng $60\text{ N}$, cùng tác dụng vào một vật tại điểm $O$. Biết góc giữa hai vectơ lực bằng $120^\circ$. Tính độ lớn của hợp lực $\vec{F} = \vec{F}_1 + \vec{F}_2$ tác dụng lên vật (theo đơn vị Niutơn).\answerbox
\oli{9.0cm}

\newpage

\item Một máy bay chuyển động đều về hướng Đông với vận tốc riêng là $600\text{ km/h}$. Cùng lúc đó có một luồng gió thổi từ hướng Bắc về hướng Nam với vận tốc $800\text{ km/h}$. Tính tốc độ thực tế của máy bay so với mặt đất (theo đơn vị $\text{km/h}$).\answerbox
\oli{9.0cm}

\item Một chiếc ca nô chuyển động thẳng đều vượt qua một con sông có hai bờ song song với vận tốc đối với dòng nước là $12\text{ km/h}$ theo phương vuông góc với bờ. Vận tốc của dòng nước chảy xuôi dòng là $5\text{ km/h}$. Tính độ lớn vận tốc của ca nô đối với bờ sông (theo đơn vị $\text{km/h}$).\answerbox
\oli{9.0cm}

\newpage

\item Cho tam giác đều $ABC$ có cạnh bằng $6\text{ cm}$. Gọi $M$ là trung điểm của cạnh $BC$. Tính độ dài của vectơ $\vec{u} = \vec{AB} - \vec{AC}$ (theo đơn vị xentimét).\answerbox
\oli{9.0cm}

\item Cho hình bình hành $ABCD$ có tâm $O$. Biết rằng độ dài của vectơ tổng $|\vec{AB} + \vec{AD}| = 14$ và độ dài của vectơ hiệu $|\vec{AB} - \vec{AD}| = 10$. Tính độ dài của đường chéo $AC$.\answerbox
\oli{9.0cm}

\end{enumerate}

\newpage

\parttitle{CHỦ ĐỀ 3. TÍCH CỦA MỘT VECTƠ VỚI MỘT SỐ (Bài 9 — Câu 21 đến Câu 30)}

\begin{enumerate}[resume*]

\item Cho đoạn thẳng $AB$. Điểm $M$ nằm giữa hai điểm $A$ và $B$ sao cho $MA = 3MB$. Biết $\vec{MA} = k\vec{MB}$ với $k$ là một số thực. Tìm giá trị của $k$.\answerbox
\oli{8.5cm}

\item Cho tam giác $ABC$ có trọng tâm $G$. Gọi $M$ là trung điểm của cạnh $BC$. Biết vectơ $\vec{AM} = k\vec{AG}$ với $k$ là số thực. Tìm giá trị của $k$.\answerbox
\oli{8.5cm}

\newpage

\item Cho tam giác $ABC$. Gọi $M$ là trung điểm của cạnh $AB$ và điểm $N$ thuộc cạnh $AC$ sao cho $AN = \dfrac{1}{3}AC$. Gọi $K$ là trung điểm của đoạn thẳng $MN$. Khi phân tích vectơ $\vec{AK}$ theo hai vectơ $\vec{AB}$ và $\vec{AC}$ dưới dạng $\vec{AK} = x\vec{AB} + y\vec{AC}$, hãy tính giá trị biểu thức $P = 12x + 18y$.\answerbox
\oli{9.0cm}

\item Cho tam giác $ABC$. Lấy điểm $D$ thuộc cạnh $BC$ sao cho $BD = 3DC$. Phân tích vectơ $\vec{AD}$ theo hai vectơ $\vec{AB}$ và $\vec{AC}$ dưới dạng $\vec{AD} = m\vec{AB} + n\vec{AC}$. Tính giá trị của biểu thức $S = 4m + 8n$.\answerbox
\oli{9.0cm}

\newpage

\item Cho đoạn thẳng $AB$. Điểm $I$ thỏa mãn hệ thức vectơ $\vec{IA} + 4\vec{IB} = \vec{0}$. Với mọi điểm $M$ bất kì, ta luôn có $\vec{MA} + 4\vec{MB} = k\vec{MI}$. Tìm giá trị của số thực $k$.\answerbox
\oli{9.0cm}

\item Cho tam giác $ABC$ có trọng tâm $G$. Gọi $I$ là trung điểm của đoạn thẳng $AG$. Phân tích vectơ $\vec{BI}$ theo hai vectơ $\vec{BA}$ và $\vec{BC}$ dưới dạng $\vec{BI} = a\vec{BA} + b\vec{BC}$. Tính giá trị biểu thức $T = 6a + 6b$.\answerbox
\oli{9.0cm}

\newpage

\item Cho tam giác đều $ABC$ có cạnh bằng $9\text{ cm}$ và trọng tâm $G$. Tập hợp các điểm $M$ thỏa mãn đẳng thức $|\vec{MA} + \vec{MB} + \vec{MC}| = 18\text{ cm}$ là một đường tròn có bán kính $R$. Tính bán kính $R$ (theo đơn vị xentimét).\answerbox
\oli{9.0cm}

\item Cho hình chữ nhật $ABCD$ có $AB = 6\text{ cm}$ và $BC = 8\text{ cm}$. Gọi $O$ là giao điểm của hai đường chéo $AC$ và $BD$. Tập hợp các điểm $M$ thỏa mãn đẳng thức $|\vec{MA} + \vec{MB} + \vec{MC} + \vec{MD}| = 20\text{ cm}$ là một đường tròn bán kính $R$. Tính bán kính $R$ (theo đơn vị xentimét).\answerbox
\oli{9.0cm}

\newpage

\item \noindent\begin{minipage}[t]{0.65\linewidth}
Một bức tranh có khối lượng $m = 4\text{ kg}$ (trọng lượng $P = 40\text{ N}$) được treo cân bằng trên tường nhờ hai sợi dây nhẹ không dãn có cùng chiều dài gắn vào đinh treo. Hai sợi dây hợp với nhau một góc bằng $120^\circ$ và đối xứng qua phương thẳng đứng (như hình vẽ bên). Tính độ lớn lực căng $T$ của mỗi sợi dây (theo đơn vị Niutơn).\answerbox
\end{minipage}\hfill
\begin{minipage}[t]{0.31\linewidth}\raggedleft
\begin{tikzpicture}[scale=0.55, >=Stealth]
  \coordinate (O) at (0,3.5);
  \coordinate (A) at (-2,0);
  \coordinate (B) at (2,0);
  \draw[thick, fill=gray!15] (-2.5,-2.2) rectangle (2.5,0);
  \draw[line width=1.1pt, accent] (O) -- (A) node[midway, left, font=\scriptsize] {$\vec{T}_1$};
  \draw[line width=1.1pt, accent] (O) -- (B) node[midway, right, font=\scriptsize] {$\vec{T}_2$};
  \draw[dashed, gray!80] (O) -- (0,-0.5);
  \fill[red!75!black] (O) circle (2.5pt) node[above, font=\scriptsize] {Đinh};
  \fill[ink] (A) circle (2pt);
  \fill[ink] (B) circle (2pt);
  \node[font=\scriptsize] at (0,-1.1) {Tranh ($4\text{ kg}$)};
  \draw[->, line width=1.2pt, red!75!black] (0,-1.1) -- (0,-2.6) node[right, font=\scriptsize] {$\vec{P}$};
  \draw[gray!70] (-0.8,2.3) arc[start angle=240, end angle=300, radius=1.4];
  \node[font=\scriptsize, gray!90!black] at (0,2.1) {$120^\circ$};
\end{tikzpicture}
\end{minipage}
\oli{7.5cm}

\item Cho tam giác $ABC$. Gọi $M$ là trung điểm của cạnh $BC$ và $N$ là trung điểm của đoạn thẳng $AM$. Đường thẳng $BN$ cắt cạnh $AC$ tại điểm $P$. Biết rằng $\vec{AC} = k\vec{AP}$ với $k$ là một số thực. Tìm giá trị của $k$.\answerbox
\oli{8.5cm}

\end{enumerate}

\newpage

\parttitle{CHỦ ĐỀ 4. VECTƠ TRONG MẶT PHẲNG TỌA ĐỘ (Bài 10 — Câu 31 đến Câu 40)}

\begin{enumerate}[resume*]

\item Trong mặt phẳng tọa độ $Oxy$, cho tam giác $ABC$ có $A(1; 4)$, $B(3; -2)$ và $C(5; 7)$. Gọi $G(x_G; y_G)$ là tọa độ trọng tâm của tam giác $ABC$. Tính giá trị của tổng $x_G + y_G$.\answerbox
\oli{8.5cm}

\item Trong mặt phẳng tọa độ $Oxy$, cho hai điểm $A(-1; 2)$ và $B(7; 6)$. Gọi $M$ là trung điểm của đoạn thẳng $AB$. Tính khoảng cách $OM$ từ điểm $M$ đến gốc tọa độ $O$.\answerbox
\oli{8.5cm}

\newpage

\item Trong mặt phẳng tọa độ $Oxy$, cho ba điểm $A(1; 2)$, $B(4; 6)$ và $C(-2; 3)$. Tìm tọa độ đỉnh $D(x_D; y_D)$ sao cho tứ giác $ABCD$ là hình bình hành. Tính giá trị của biểu thức $x_D + 2y_D$.\answerbox
\oli{9.0cm}

\item Trong mặt phẳng tọa độ $Oxy$, cho hai vectơ $\vec{u} = (2; -3)$ và $\vec{v} = (1; 4)$. Vectơ $\vec{w} = 3\vec{u} - 2\vec{v}$ có tọa độ là $(x_0; y_0)$. Tính giá trị của hiệu $x_0 - y_0$.\answerbox
\oli{9.0cm}

\newpage

\item Trong mặt phẳng tọa độ $Oxy$, cho ba điểm $A(2; 1)$, $B(5; 7)$ và $C(m; 3)$ với $m$ là tham số thực. Tìm giá trị của $m$ để ba điểm $A, B, C$ thẳng hàng.\answerbox
\oli{9.0cm}

\item Trong mặt phẳng tọa độ $Oxy$, cho hai vectơ $\vec{a} = (2; 1)$ và $\vec{b} = (-1; 3)$. Biết rằng vectơ $\vec{c} = (7; 7)$ được biểu diễn dưới dạng $\vec{c} = x\vec{a} + y\vec{b}$. Tính giá trị của $x^2 + y^2$.\answerbox
\oli{9.0cm}

\newpage

\item Trong mặt phẳng tọa độ $Oxy$, cho hai điểm $A(1; 3)$ và $B(5; -1)$. Điểm $M(x_M; 0)$ thuộc trục hoành $Ox$ sao cho tam giác $MAB$ cân tại đỉnh $M$. Tính giá trị của $x_M$.\answerbox
\oli{9.0cm}

\item Trong mặt phẳng tọa độ $Oxy$, cho tam giác $ABC$ có $A(-2; 5)$, $B(1; 2)$ và $C(7; 10)$. Gọi $M, N$ lần lượt là trung điểm của các cạnh $AB$ và $AC$. Tính độ dài đoạn thẳng $MN$.\answerbox
\oli{9.0cm}

\newpage

\item Một ca nô chuyển động thẳng đều trên mặt biển gắn với hệ trục tọa độ $Oxy$ (đơn vị trên mỗi trục là kilômét). Lúc $7$ giờ sáng ca nô ở vị trí $A(3; 2)$, đến $9$ giờ sáng ca nô ở vị trí $B(9; 10)$. Nếu giữ nguyên vận tốc và hướng đi, đến $10$ giờ sáng cùng ngày ca nô ở vị trí $C(x_C; y_C)$. Tính giá trị của hoành độ $x_C$.\answerbox
\oli{9.0cm}

\item Trong mặt phẳng tọa độ $Oxy$, cho ba điểm $A(1; 2)$, $B(3; 5)$ và $C(6; 1)$. Điểm $M(x_M; y_M)$ thỏa mãn đẳng thức vectơ $\vec{MA} + 2\vec{MB} - \vec{MC} = \vec{0}$. Tính giá trị của tổng $x_M + y_M$.\answerbox
\oli{9.0cm}

\end{enumerate}


\newpage
\pagestyle{dapan}
\setcounter{page}{1}
\begin{center}
  {\Large\bfseries\color{ink} BẢNG ĐÁP ÁN \& HƯỚNG DẪN GIẢI CHI TIẾT}\\[3pt]
  {\small\color{gray!75!black} Đề cương ôn tập Toán 10 — Chương IV: Vectơ (Bài 7, 8, 9, 10 — 40 câu trả lời ngắn)}
\end{center}
\vspace{2pt}


\begin{center}
\renewcommand{\arraystretch}{1.12}
\begin{tabular}{|C{1.1cm}|C{2.0cm}|C{1.1cm}|C{2.0cm}|C{1.1cm}|C{2.0cm}|C{1.1cm}|C{2.0cm}|}
\hline
\textbf{Câu} & \textbf{Đáp án} & \textbf{Câu} & \textbf{Đáp án} & \textbf{Câu} & \textbf{Đáp án} & \textbf{Câu} & \textbf{Đáp án} \\ \hline
\textbf{1} & 90 & \textbf{11} & 13 & \textbf{21} & -3 & \textbf{31} & 6 \\ \hline
\textbf{2} & 4 & \textbf{12} & 16 & \textbf{22} & 1,5 & \textbf{32} & 5 \\ \hline
\textbf{3} & 2 & \textbf{13} & 9 & \textbf{23} & 6 & \textbf{33} & -7 \\ \hline
\textbf{4} & 5 & \textbf{14} & 10 & \textbf{24} & 7 & \textbf{34} & 21 \\ \hline
\textbf{5} & 5 & \textbf{15} & 50 & \textbf{25} & 5 & \textbf{35} & 3 \\ \hline
\textbf{6} & 7,5 & \textbf{16} & 60 & \textbf{26} & 5 & \textbf{36} & 17 \\ \hline
\textbf{7} & 32 & \textbf{17} & 1000 & \textbf{27} & 6 & \textbf{37} & 2 \\ \hline
\textbf{8} & 4 & \textbf{18} & 13 & \textbf{28} & 5 & \textbf{38} & 5 \\ \hline
\textbf{9} & 10 & \textbf{19} & 6 & \textbf{29} & 40 & \textbf{39} & 12 \\ \hline
\textbf{10} & 5 & \textbf{20} & 14 & \textbf{30} & 3 & \textbf{40} & 6 \\ \hline
\end{tabular}
\end{center}

\vspace{2pt}

\setlength{\columnsep}{16pt}

\parttitle{HƯỚNG DẪN GIẢI CHI TIẾT (CÂU 1 – CÂU 40)}

\begin{multicols}{2}

\small\sloppy

\needspace{2\baselineskip}\textbf{Câu 1.} Mỗi cách chọn $2$ đỉnh có thứ tự từ $10$ đỉnh tạo thành một vectơ khác $\vec{0}$. Số vectơ là $A_{10}^2 = 10 \cdot 9 = \mathbf{90}$.\par\vspace{2.5pt}

\needspace{2\baselineskip}\textbf{Câu 2.} Đường thẳng $AB$ chứa $A, B \Rightarrow \vec{AB}, \vec{BA}$ ($2$ vectơ). Đường thẳng $CD$ chứa $C, D \Rightarrow \vec{CD}, \vec{DC}$ ($2$ vectơ). Tâm $O$ không nằm trên $AB, CD$. Vậy có $2 + 2 = \mathbf{4}$ vectơ.\par\vspace{2.5pt}

\needspace{2\baselineskip}\textbf{Câu 3.} Do $MN$ là đường trung bình của $\triangle ABC$ nên $\vec{MN} = \vec{BP} = \vec{PA}$. Vậy có $\mathbf{2}$ vectơ khác $\vec{MN}$ bằng nó.\par\vspace{2.5pt}

\needspace{2\baselineskip}\textbf{Câu 4.} Vectơ $\vec{AB}$ có độ dời ngang bằng $3$ và độ dời dọc bằng $4$. Độ dài $|\vec{u}| = \sqrt{3^2 + 4^2} = \mathbf{5}$.\par\vspace{2.5pt}

\needspace{2\baselineskip}\textbf{Câu 5.} Kẻ $BH \perp CD$ tại $H \Rightarrow DH = AB = 3\text{ cm} \Rightarrow HC = 7 - 3 = 4\text{ cm}$. Tam giác vuông $BHC$ có $BH = AD = 3\text{ cm} \Rightarrow BC = \sqrt{3^2 + 4^2} = \mathbf{5}\text{ cm}$.\par\vspace{2.5pt}

\needspace{2\baselineskip}\textbf{Câu 6.} Cạnh huyền $BC = \sqrt{9^2 + 12^2} = 15\text{ cm}$. Trung tuyến ứng với cạnh huyền $AM = \dfrac{1}{2}BC = \mathbf{7{,}5}\text{ cm}$.\par\vspace{2.5pt}

\needspace{2\baselineskip}\textbf{Câu 7.} Tam giác $ABC$ vuông cân có $AC = 8\sqrt{2}\text{ cm} \Rightarrow MN = \dfrac{1}{2}AC = 4\sqrt{2}\text{ cm}$. Do đó $|\vec{MN}|^2 = (4\sqrt{2})^2 = \mathbf{32}$.\par\vspace{2.5pt}

\needspace{2\baselineskip}\textbf{Câu 8.} $DE$ là đường trung bình của $\triangle ABC$ ứng với đáy $BC \Rightarrow DE = \dfrac{1}{2}BC = \dfrac{1}{2}\cdot 8 = \mathbf{4}$.\par\vspace{2.5pt}

\needspace{2\baselineskip}\textbf{Câu 9.} Cạnh $AD = \dfrac{48}{6} = 8\text{ cm}$. Đường chéo $BD = \sqrt{AB^2 + AD^2} = \sqrt{6^2 + 8^2} = \mathbf{10}\text{ cm}$.\par\vspace{2.5pt}

\needspace{2\baselineskip}\textbf{Câu 10.} Với điểm đầu cố định là $A$, có $5$ cách chọn điểm cuối trong các điểm còn lại. Vậy có $\mathbf{5}$ vectơ.\par\vspace{2.5pt}

\needspace{2\baselineskip}\textbf{Câu 11.} $\vec{u} = \vec{AB} - \vec{CB} = \vec{AB} + \vec{BC} = \vec{AC} \Rightarrow |\vec{u}| = AC = \sqrt{5^2 + 12^2} = \mathbf{13}$.\par\vspace{2.5pt}

\needspace{2\baselineskip}\textbf{Câu 12.} Theo quy tắc hình bình hành: $\vec{AB} + \vec{AD} = \vec{AC}$. Giao điểm hai đường chéo $O$ cho $OB = 6\text{ cm} \Rightarrow OA = \sqrt{10^2 - 6^2} = 8\text{ cm} \Rightarrow AC = 16\text{ cm}$. Vậy $|\vec{w}| = \mathbf{16}$.\par\vspace{2.5pt}

\needspace{2\baselineskip}\textbf{Câu 13.} $\vec{x} = (\vec{AB} - \vec{AD}) + \vec{CD} = \vec{DB} + \vec{CD} = \vec{CD} + \vec{DB} = \vec{CB} \Rightarrow |\vec{x}| = BC = \mathbf{9}\text{ cm}$.\par\vspace{2.5pt}

\needspace{2\baselineskip}\textbf{Câu 14.} $\vec{y} = \vec{AB} - \vec{AD} = \vec{DB} \Rightarrow |\vec{y}| = BD = \sqrt{8^2 + 6^2} = \mathbf{10}$.\par\vspace{2.5pt}

\needspace{2\baselineskip}\textbf{Câu 15.} Hai lực vuông góc nên độ lớn hợp lực là $F = \sqrt{F_1^2 + F_2^2} = \sqrt{30^2 + 40^2} = \mathbf{50}\text{ N}$.\par\vspace{2.5pt}

\needspace{2\baselineskip}\textbf{Câu 16.} $F = \sqrt{60^2 + 60^2 + 2\cdot 60\cdot 60\cdot \cos 120^\circ} = \sqrt{3600 + 3600 - 3600} = \mathbf{60}\text{ N}$.\par\vspace{2.5pt}

\needspace{2\baselineskip}\textbf{Câu 17.} Vận tốc máy bay và vận tốc gió vuông góc nhau, vận tốc thực tế $v = \sqrt{600^2 + 800^2} = \mathbf{1000}\text{ km/h}$.\par\vspace{2.5pt}

\needspace{2\baselineskip}\textbf{Câu 18.} Vận tốc ca nô đối với bờ là hợp của hai vận tốc vuông góc: $v = \sqrt{12^2 + 5^2} = \mathbf{13}\text{ km/h}$.\par\vspace{2.5pt}

\needspace{2\baselineskip}\textbf{Câu 19.} $\vec{u} = \vec{AB} - \vec{AC} = \vec{CB} \Rightarrow |\vec{u}| = CB = \mathbf{6}\text{ cm}$.\par\vspace{2.5pt}

\needspace{2\baselineskip}\textbf{Câu 20.} $\vec{AB} + \vec{AD} = \vec{AC} \Rightarrow AC = |\vec{AB} + \vec{AD}| = \mathbf{14}$.\par\vspace{2.5pt}

\needspace{2\baselineskip}\textbf{Câu 21.} Vì $M$ nằm giữa $A$ và $B$ nên $\vec{MA}$ và $\vec{MB}$ ngược hướng. Từ $MA = 3MB \Rightarrow \vec{MA} = -3\vec{MB} \Rightarrow k = \mathbf{-3}$.\par\vspace{2.5pt}

\needspace{2\baselineskip}\textbf{Câu 22.} $G$ là trọng tâm $\Rightarrow AG = \dfrac{2}{3}AM \Rightarrow AM = \dfrac{3}{2}AG \Rightarrow \vec{AM} = 1{,}5\vec{AG} \Rightarrow k = \mathbf{1{,}5}$.\par\vspace{2.5pt}

\needspace{2\baselineskip}\textbf{Câu 23.} $\vec{AK} = \dfrac{1}{2}(\vec{AM} + \vec{AN}) = \dfrac{1}{4}\vec{AB} + \dfrac{1}{6}\vec{AC} \Rightarrow x = \dfrac{1}{4},\; y = \dfrac{1}{6} \Rightarrow P = 12\left(\dfrac{1}{4}\right) + 18\left(\dfrac{1}{6}\right) = \mathbf{6}$.\par\vspace{2.5pt}

\needspace{2\baselineskip}\textbf{Câu 24.} Ta có $\vec{AD} = \dfrac{1}{4}\vec{AB} + \dfrac{3}{4}\vec{AC} \Rightarrow m = \dfrac{1}{4},\; n = \dfrac{3}{4}$. Vậy $S = 4\left(\dfrac{1}{4}\right) + 8\left(\dfrac{3}{4}\right) = 1 + 6 = \mathbf{7}$.\par\vspace{2.5pt}

\needspace{2\baselineskip}\textbf{Câu 25.} $\vec{MA} + 4\vec{MB} = 5\vec{MI} + (\vec{IA} + 4\vec{IB}) = 5\vec{MI} \Rightarrow k = \mathbf{5}$.\par\vspace{2.5pt}

\needspace{2\baselineskip}\textbf{Câu 26.} $\vec{BG} = \dfrac{1}{3}\vec{BA} + \dfrac{1}{3}\vec{BC}$. Do $I$ là trung điểm $AG$ nên $\vec{BI} = \dfrac{1}{2}(\vec{BA} + \vec{BG}) = \dfrac{2}{3}\vec{BA} + \dfrac{1}{6}\vec{BC} \Rightarrow a = \dfrac{2}{3},\; b = \dfrac{1}{6} \Rightarrow T = 6\left(\dfrac{2}{3}\right) + 6\left(\dfrac{1}{6}\right) = 4 + 1 = \mathbf{5}$.\par\vspace{2.5pt}

\needspace{2\baselineskip}\textbf{Câu 27.} Ta có $|\vec{MA} + \vec{MB} + \vec{MC}| = |3\vec{MG}| = 3MG = 18\text{ cm} \Leftrightarrow MG = 6\text{ cm}$. Vậy bán kính đường tròn $R = \mathbf{6}\text{ cm}$.\par\vspace{2.5pt}

\needspace{2\baselineskip}\textbf{Câu 28.} $|\vec{MA} + \vec{MB} + \vec{MC} + \vec{MD}| = |4\vec{MO}| = 4MO = 20\text{ cm} \Leftrightarrow MO = 5\text{ cm}$. Vậy bán kính đường tròn $R = \mathbf{5}\text{ cm}$.\par\vspace{2.5pt}

\needspace{2\baselineskip}\textbf{Câu 29.} Gọi góc hợp bởi mỗi dây với phương thẳng đứng là $60^\circ$. Điều kiện cân bằng lực: $2T\cos 60^\circ = P \Leftrightarrow 2T \cdot 0{,}5 = 40 \Leftrightarrow T = \mathbf{40}\text{ N}$.\par\vspace{2.5pt}

\needspace{2\baselineskip}\textbf{Câu 30.} Kẻ $MK \parallel BP$ ($K \in AC$) được $P$ là trung điểm của $CK$ và $AK$. Do đó $AP = PK = KC \Rightarrow AC = 3AP \Rightarrow k = \mathbf{3}$.\par\vspace{2.5pt}

\needspace{2\baselineskip}\textbf{Câu 31.} $x_G = \dfrac{1 + 3 + 5}{3} = 3$, $y_G = \dfrac{4 - 2 + 7}{3} = 3 \Rightarrow x_G + y_G = 3 + 3 = \mathbf{6}$.\par\vspace{2.5pt}

\needspace{2\baselineskip}\textbf{Câu 32.} Trung điểm $M$ có tọa độ $M(3; 4) \Rightarrow OM = \sqrt{3^2 + 4^2} = \mathbf{5}$.\par\vspace{2.5pt}

\needspace{2\baselineskip}\textbf{Câu 33.} $ABCD$ là hình bình hành $\Leftrightarrow \vec{AD} = \vec{BC} \Leftrightarrow (x_D - 1; y_D - 2) = (-6; -3) \Leftrightarrow D(-5; -1)$. Vậy $x_D + 2y_D = -5 + 2(-1) = \mathbf{-7}$.\par\vspace{2.5pt}

\needspace{2\baselineskip}\textbf{Câu 34.} $\vec{w} = (3\cdot 2 - 2\cdot 1;\; 3(-3) - 2\cdot 4) = (4; -17) \Rightarrow x_0 - y_0 = 4 - (-17) = \mathbf{21}$.\par\vspace{2.5pt}

\needspace{2\baselineskip}\textbf{Câu 35.} $\vec{AB} = (3; 6)$, $\vec{AC} = (m - 2; 2)$. Ba điểm thẳng hàng $\Leftrightarrow \dfrac{m - 2}{3} = \dfrac{2}{6} = \dfrac{1}{3} \Leftrightarrow m - 2 = 1 \Leftrightarrow m = \mathbf{3}$.\par\vspace{2.5pt}

\needspace{2\baselineskip}\textbf{Câu 36.} Ta có hệ $2x - y = 7$ và $x + 3y = 7 \Leftrightarrow x = 4,\; y = 1 \Rightarrow x^2 + y^2 = 4^2 + 1^2 = \mathbf{17}$.\par\vspace{2.5pt}

\needspace{2\baselineskip}\textbf{Câu 37.} $MA^2 = MB^2 \Leftrightarrow (x_M - 1)^2 + 9 = (x_M - 5)^2 + 1 \Leftrightarrow 8x_M = 16 \Leftrightarrow x_M = \mathbf{2}$.\par\vspace{2.5pt}

\needspace{2\baselineskip}\textbf{Câu 38.} $MN$ là đường trung bình của $\triangle ABC \Rightarrow MN = \dfrac{1}{2}BC = \dfrac{1}{2}\sqrt{(7 - 1)^2 + (10 - 2)^2} = \dfrac{1}{2}\sqrt{100} = \mathbf{5}$.\par\vspace{2.5pt}

\needspace{2\baselineskip}\textbf{Câu 39.} Trong $2\text{ h}$ ca nô đi vectơ $\vec{AB} = (6; 8) \Rightarrow$ vận tốc $\vec{v} = (3; 4)\text{ km/h}$. Thêm $1\text{ h}$ đến $10\text{h}$, tọa độ $C = B + \vec{v} = (9 + 3; 10 + 4) = (12; 14) \Rightarrow x_C = \mathbf{12}$.\par\vspace{2.5pt}

\needspace{2\baselineskip}\textbf{Câu 40.} $\vec{MA} + 2\vec{MB} - \vec{MC} = \vec{0} \Leftrightarrow \vec{OM} = \dfrac{\vec{OA} + 2\vec{OB} - \vec{OC}}{2}$. Tọa độ $x_M = \dfrac{1 + 6 - 6}{2} = 0{,}5$; $y_M = \dfrac{2 + 10 - 1}{2} = 5{,}5 \Rightarrow x_M + y_M = \mathbf{6}$.\par\vspace{2.5pt}

\end{multicols}

\end{document}